% Generated by scripts/assemble.py; corrections belong in source/corrections.json.
\examhead{101}{浙江大学 2024 年数学分析试卷}{100}
% p101-b20; source PDF page 101
\blockanchor{p101-b20}
\question{1}计算题.\par
% p101-b21; source PDF page 101
\par\Needspace{6\baselineskip}
\blockanchor{p101-b21}
(1) 求极限\par
\nopagebreak[4]
% p101-b22; source PDF page 101
\blockanchor{p101-b22}
\displayformula{\lim _ {n \rightarrow \infty} \sqrt [ n ]{\cos \frac {\pi}{2 n} \cdot \cos \frac {2 \pi}{2 n} \cdot \cos \frac {(n - 1) \pi}{2 n}}.}
% p101-b23; source PDF page 101
\par\Needspace{6\baselineskip}
\blockanchor{p101-b23}
(2) 求积分\par
\nopagebreak[4]
% p101-b24; source PDF page 101
\blockanchor{p101-b24}
\displayformula{\iint_ {x ^ {2} + y ^ {2} \leq 1} | 3 x + 4 y | \mathrm{d} x \mathrm{d} y.}
% p101-b25; source PDF page 101
\par\Needspace{6\baselineskip}
\blockanchor{p101-b25}
(3) 求极限\par
\nopagebreak[4]
% p101-b26; source PDF page 101
\blockanchor{p101-b26}
\displayformula{\lim _ {x \rightarrow 0} \left[ \frac {1}{\ln (1 + x)} - \frac {1}{\ln (x + \sqrt {1 + x ^ {2}})} \right].}
% p102-b00; source PDF page 102
\par\Needspace{6\baselineskip}
\blockanchor{p102-b00}
(4) 求极限\par
\nopagebreak[4]
% p102-b01; source PDF page 102
\blockanchor{p102-b01}
\displayformula{\lim _ {x \rightarrow + \infty} \frac {1}{x} \int_ {0} ^ {x} | \sin t | \mathrm{d} t.}
% p102-b02; source PDF page 102
\blockanchor{p102-b02}
\question{2}解答如下问题:\par
% p102-b03; source PDF page 102
\blockanchor{p102-b03}
(1) 请叙述确界存在定理.\par
% p102-b04; source PDF page 102
\blockanchor{p102-b04}
(2) 请利用闭区间套定理证明确界存在定理.\par
% p102-b05; source PDF page 102
\blockanchor{p102-b05}
\question{3}设  $x_{1}=1, x_{n+1}=\sqrt{\frac{2x_{n}^{2}}{2+x_{n}^{2}}}$ ,  $n=1,2,3,\cdots$\par
% p102-b06; source PDF page 102
\blockanchor{p102-b06}
(1) 证明: $\lim_{n\to \infty}x_n = 0.$\par
% p102-b07; source PDF page 102
\blockanchor{p102-b07}
(2) 计算 $\lim_{n\to \infty}n(x_n - \ln (1 + x_n))$\par
% p102-b08; source PDF page 102
\par\Needspace{6\baselineskip}
\blockanchor{p102-b08}
\question{4}设函数 $f_{1}(x)$ 在 $[0,1]$ 上连续, $K(x,y)$ 在 $[0,1] \times [0,1]$ 上连续. 对任意 $x \in [0,1]$, 定义\par
\nopagebreak[4]
% p102-b09; source PDF page 102
\blockanchor{p102-b09}
\displayformula{f _ {n + 1} (x) = \int_ {0} ^ {x} K (x, y) f _ {n} (y) \mathrm{d} y, n = 1, 2, \dots .}
% p102-b10; source PDF page 102
\blockanchor{p102-b10}
请证明：函数列 $\{f_n(x)\}$ 在[0,1]上一致收敛于0.\par
% p102-b11; source PDF page 102
\blockanchor{p102-b11}
\question{5}设函数 $f, g$ 在 $[0,1]$ 上连续, 且存在包含于 $[0,1]$ 的数列 $\{x_n\}$, 使得对于任意 $n \geq 1$, 有 $f(x_n) = g(x_{n+1})$. 请证明: 存在 $\xi \in [0,1]$, 使得 $f(\xi) = g(\xi)$.\par
% p102-b12; source PDF page 102
\blockanchor{p102-b12}
\question{6}设 $f$ 是定义在 $[a, b]$ 上的函数, 且 $G = \{(x, f(x)) \mid x \in [a, b]\}$ 是 $\mathbb{R}^2$ 上的有界闭集. 请证明: $f$ 是 $[a, b]$ 上的连续函数.\par
% p102-b13; source PDF page 102
\par\Needspace{6\baselineskip}
\blockanchor{p102-b13}
\question{7}设函数 $f$ 是 (0,1] 上的无界的单调函数, 且广义积分 $\int_{0}^{1} f(x) \, \mathrm{d}x$ 收敛. 证明:\par
\nopagebreak[4]
% p102-b14; source PDF page 102
\blockanchor{p102-b14}
\displayformula{\lim _ {n \rightarrow \infty} \frac {1}{n} \sum_ {k = 2} ^ {n} f \left(\frac {k - 1}{n}\right) = \int_ {0} ^ {1} f (x) \mathrm{d} x.}
% p102-b15; source PDF page 102
\blockanchor{p102-b15}
\question{8}设函数 $f$ 在 $\mathbb{R}$ 上可微, 且满足对任意 $x \in \mathbb{R}$, 有 $f(x + 1) - f(x) = f'(x)$ 以及 $\lim_{x \to +\infty} f'(x) = 1$. 请证明: 存在常数 $C$, 使得 $f(x) = x + C$.\par
% p102-b16; source PDF page 102
\par\Needspace{6\baselineskip}
\blockanchor{p102-b16}
\question{9}设函数 $f$ 在 $[0,1]$ 上连续, 定义\par
\nopagebreak[4]
% p102-b17; source PDF page 102
\blockanchor{p102-b17}
\displayformula{g (t) = \int_ {0} ^ {1} \frac {t f (x)}{x ^ {2} + t ^ {2}} \mathrm{d} x, t \in \mathbb {R}.}
% p102-b18; source PDF page 102
\blockanchor{p102-b18}
请证明: 函数 g 在点 0 处连续当且仅当  $f(0)=0$ .\par
